\documentclass[10pt,a4paper]{amsart}
\usepackage[margin=19mm]{geometry}
\usepackage{amsmath,amssymb,mathtools}
\usepackage[scr=rsfs,cal=euler]{mathalfa}
\usepackage{xfrac}
\usepackage[unicode,hidelinks,bookmarks=false]{hyperref}
\newcommand{\ie}{i.e.\ }
\newcommand{\cf}{cf.\ }
\newcommand{\resp}{respectively\ }
\newcommand{\Subsection}{\subsection}
\providecommand{\nobreakdash}{\nobreak\hbox{-}}
\setlength{\emergencystretch}{2em}
\multlinegap0pt
\usepackage{amssymb}
\usepackage{enumerate}
\usepackage{xcolor}
\usepackage{tikz}
\newtheorem{assumption}{Assumption}
\newtheorem{lemma}{Lemma}
\newcommand{\De}{\Delta}
\newcommand{\Ga}{\Gamma}
\newcommand{\Rn}{{\mathbb R}^{n-1}}
\newcommand {\al}{\alpha}
\newcommand {\be}{\beta}
\newcommand{\om}{{ \omega  }}
\title{On the Existence of Boundary Layer Separation for Incompressible Fluid Flow in the Half-Space}
\author{Tongkeun Chang, Kyungkeun Kang}
\date{Source: arXiv:2605.08573v1 (CC BY 4.0)}
\begin{document}
\maketitle

\noindent\emph{Source and licence.} On the Existence of Boundary Layer Separation for Incompressible Fluid Flow in the Half-Space. Version arXiv:2605.08573v1 (CC BY 4.0), distributed by arXiv under the Creative Commons Attribution 4.0 International licence, http://creativecommons.org/licenses/by/4.0/, as recorded in the arXiv licence metadata for this exact version and shown on the abstract page https://arxiv.org/abs/2605.08573v1. Full licence notice: \texttt{licenses/arxiv-2605.08573-CC-BY-4.0.txt}.\par\smallskip

\noindent\textit{Editorial scope.} Counted unit: the article's lemma theorem1221-2 (source0.tex lines 1182--1191). The article's complete original proof of it is delivered with it (source0.tex lines 1193--1294, one \texttt{proof} environment of 102 lines). No mathematics is written, repaired or re-derived: the statement and the proof are the article's own lines, in the article's own notation, and no retained line is substituted or re-flowed.  Retained uncounted as the source-local context the proof needs: lines 424--426; lines 552--554; lines 787--789; lines 791--793; lines 808--816.  Nothing was omitted from any retained span, so the package carries no omission record.  No retained line contains a citation, so the wrapper carries no bibliography and no bibliography entry.  No retained line contains a figure environment, an image-inclusion command or drawing code, so no image is delivered and the package declares no asset dependency.  Editorial formatting only, in addition: an AMS \texttt{amsart} preamble that restates the article's own declarations and macros used by the retained lines, namely the environments \texttt{assumption}, \texttt{lemma} on separate counters, together with the standard packages those lines need.  The retained environments print in this document as Assumption 1, Lemma 1 in order of appearance, whereas the article prints the same items with the numbering of its own full document.  The complete native source of the article as distributed in the arXiv e-print for this exact version is bundled unchanged as \texttt{sources/200160/source0.tex}.
\begin{assumption}\label{assume1}
Let $\phi \in C^{\frac12 +\frac{\al}2} ([0, 2 ) )$ for some $ 0 < \al  <1$,  and $ \phi$ be monotonically increasing in $(0,1)$ and monotonically decreasing in $(1,2)$ with $\phi(0) =0$.   
\end{assumption}

\begin{align}\label{260321-7}
M'(t) < -e_0.
\end{align}

 \begin{align} \label{shortg}
D_{x_n} w_i (x', 0,t) & =  - \frac1{\pi^\frac12}  \int_0^t  \frac{ \phi(s)}{(t-s)^\frac32}     f_i (x', t-s) ds    + \frac1{\pi^\frac12}  R'_i \psi(x) \int_{-\infty}^t     \frac{\phi(t) - \phi(s)}{(t-s)^\frac32}  ds, \quad  1 \leq i \leq n-1, 
\end{align}

\begin{align}\label{defoff}
  f_i (x', t-s)  = \int_{\Rn}   \Ga'(x' -y', t-s) \big( R'_i \psi(y')   - R'_i \psi(x')    \big)  dy'.
\end{align}

\begin{lemma}\label{lemma1116-1}
Let $ f_i$ , $i\neq n$  be the function defined in \eqref{defoff} and $ \be = (\be_1,\be_2, \cdots, \be_{n-1})   \in ( {\mathbb N} \cup \{ 0\} )^{n-1}$ with $ |\be| = \sum_{1 \leq k \leq n-1} \be_k$.
Then, there exists  $ c_* > 2$ such that for $ 0 < t <2$ and $ |x'| \geq c_*$, 
$D_{x'}^\be  f_i$ can be  decomposed as $ D_{x'}^\be  f_i(x',t) = F_1(x', t) + F_2(x', t)$, where 
\begin{align*}
F_1(x',t) &  =   \De'  D_{x'}^\be   R'_i \psi (x'  )    t, \qquad |F_2 (x',t)|   \leq c t^\frac32 |x'|^{-n -2 -|\be|}.
\end{align*}

\end{lemma}

\begin{lemma}\label{theorem1221-2}
Let $\phi \in C^2([0,2])$ satisfy Assumption \ref{assume1} and $ \phi'(0) =0$. Let $M$ satisfy \eqref{260321-7}.  
\begin{itemize}
\item[(i)]
Let $ 1 \leq i \leq n-1$. There exist positive constants $ c_{*5}>2$ and $c_i>0$  such that  if $c_{*5} < |x'| < c_i x_i $,  then $ D_{x_n} w_i(x', 0,t)$ is strictly decreasing  in $ t \in (1,2)$.

\item[(ii)] There exists  a  positive constant $c_{*6}>0$ such that   if $c_{*6} < |x'|   $,  then $ D^2_{x_n} w_n(x', 0,t)$ is strictly decreasing in $t \in (1,2)$.
\end{itemize}

\end{lemma}

 \begin{proof}
Since the proofs of (i) and (ii) are similar, we only prove (ii). Using the integration by parts, we have 
\begin{align}\label{0904-1}
\begin{split}
\int_0^t   \frac{\phi(s)}{(t-s)^\frac32}     D_{x_k}  f_k (x', t-s) 
%& =    2\int_0^t  \phi(s)    \frac{d}{ds}   (t-s)^{-\frac12}   f(x', t-s)  ds\\
=  -  2\int_0^t    \Big(  \phi' (s)      D_{x_k}  f _k (x', t-s)  -  \phi (s)      \frac{d}{dt} D_{x_k}  f _k (x', t-s)  \Big)  (t-s)^{-\frac12} ,
\end{split}
\end{align}
where $f_k$ is defined in \eqref{defoff}. 
Since $ \big( D_t -\De' \big) \Ga'(y', t) =0$ for $ t> 0$ and   $\int_{\Rn} \Ga' (x' - y', t) dy' =1$ for all  $(x', t)$, we have
\begin{align}\label{0904-2}
\begin{split}
\frac{d}{dt}  D_{x_k}  f _k ( x', t)  &  =\int_{\Rn} \De'_{y'} \Ga' (x'-y', t) \Big( D_{y_k} R'_k \psi(y')  - D_{x_k} R'_k \psi(x') \Big) dy'\\  
&  =\int_{\Rn}  \Ga' (x'-y', t)  \De'_{y '} D_{y_k} R'_k \psi(y')  dy'\\ 
&  =\int_{\Rn}   \Ga' (x'-y', t) \Big(  \De'_{y'} D_{y_k}   R'_k\psi(y')  -  \De'_{x'} D_{x_k} R'_k \psi(x') \Big) dy' + \De'_{x'} D_{x_k} R'_k \psi(x')\\
& := g_k(x', t) +   \De'_{x'} D_{x_k} R'_k \psi(x'),
\end{split}
\end{align}
where 
\begin{align*}
g_k (x', t-s)  = \int_{\Rn} \Ga' (x'- y', t-s) \Big(  \De'  D_{y_k} R'_k \psi(y') -   \De' D_{x_k}  R'_k \psi(x') \Big) dy'.
\end{align*} 
From \eqref{0904-1} and \eqref{0904-2},  we have 
\begin{align*}%\label{260106-1}
\begin{split}
&\int_0^t   \frac{\phi(s)}{(t-s)^\frac32}  D_{x_k}  f _k (x', t-s) ds\\
 =& -  2\int_0^t    \Big(  \phi' (s)    D_{x_k}  f _k (x', t-s)  -  \phi (s)   \big(g_k (x', t-s) +   \De'_{x'} D_{x_k} R'_k \psi(x')    \big)  \Big)  (t-s)^{-\frac12}   ds\\
=& - 4\int_0^t    \Big(  \phi'' (s)    D_{x_k}  f _k(x', t-s)  -  2\phi'  (s)   \Big(g_k (x', t-s) +   \De'_{x'} D_{x_k} R'_k \psi(x') \Big) \\
& \quad    + \phi (s) \big(h_k (x', t-s)   + ( \De')^2 D_{x_k}  R'_k \psi(x')   \big) \Big)  (t-s)^{\frac12}   ds,
\end{split}
\end{align*}



where 
\begin{align*}
h_k (x', t-s)  = \int_{\Rn} \Ga' (x'- y', t-s) \Big(  ( \De')^2  D_{y_k} R'_k \psi(y') -  ( \De')^2  D_{x_k}  R'_k \psi(x') \Big) dy'.
\end{align*}





 


 
Reminding that $ {\rm div} \, w =0$, we have  from \eqref{shortg}
\begin{align}\label{260318-1-1}
\begin{split}
D_{x_n}^2 w_n(x',  0, t) & =  -\sum_{1 \leq k \leq n-1} D_{x_k} D_{x_n} w_k (x', 0, t)\\
 & =   c_n \sum_{1 \leq k \leq n-1}\int_0^t  \frac{ \phi(s)}{(t-s)^\frac32}    D_{x_k}  f_k (x', t-s) ds    -c_n  \sum_{1 \leq k \leq n-1} D_{x_k}  R'_k \psi(x) M(t)\\
 & = -4  c_n  \sum_{1 \leq k \leq n-1} \int_0^t    \Big(  \phi'' (s)    D_{x_k}  f _k(x', t-s)  -  2\phi'  (s)   [g_k (x', t-s) +   \De'_{x'} D_{x_k} R'_k \psi(x') ] \\
& \quad    + \phi (s) \big(h_k (x', t-s)   + ( \De')^2 D_{x_k}  R'_k \psi(x')   \big) \Big)  (t-s)^{\frac12}   ds\\
 &  \quad    -4c_n  \sum_{1 \leq k \leq n-1} D_{x_k}  R'_k \psi(x) M(t).
 \end{split}
\end{align}
Hence,  it follows from \eqref{260318-1-1} that
\begin{align}\label{260318-1}
\begin{split}
\frac{d}{dt} D_{x_n}^2 w_n(x',  0, t) 
 & =  -4  c_n  \sum_{1 \leq k \leq n-1} \int_0^t    \Big(  \phi'' (s)     \frac{d}{dt}  D_{x_k}  f _k(x', t-s)  -  2\phi'  (s)   \frac{d}{dt}  g_k (x', t-s)  \\
& \quad    + \phi (s)  \frac{d}{dt}  h_k (x', t-s)   \Big)  (t-s)^{\frac12}   ds\\
&\quad   -2    c_n  \sum_{1 \leq k \leq n-1} \int_0^t    \Big(  \phi'' (s)    D_{x_k}  f _k(x', t-s)  -  2\phi'  (s)   \Big(g_k (x', t-s) +   \De'_{x'} D_{x_k} R'_k \psi(x') \Big) \\
& \quad    + \phi (s) \big(h_k (x', t-s)   + ( \De')^2 D_{x_k}  R'_k \psi(x')   \big) \Big)  (t-s)^{-\frac12}   ds\\
 &  \quad    -2c_n  \sum_{1 \leq k \leq n-1} D_{x_k}  R'_k \psi(x)M'(t).
 \end{split}
\end{align}
 
 

Note that  for  $  2 < | x'| $.
\begin{align}\label{260318-3}
\sum_{1 \leq k \leq n-1} D_{x_k}   R'_k \psi(x')  & = -    \frac1{n\om_n } \int_{\Rn}  \frac{1} {|x' -z'|^n}   \psi(z')   dz'  \approx - |x'|^{-n} c_\psi.
\end{align}
From the proof of Lemma \ref{lemma1116-1},  for $ 2 < |x'| $,
\begin{align}\label{260404-6}
\begin{split}
|     D_{x_k} f_k ( x', t)|, \quad  |x'|^2|     g_k ( x', t)|, \quad |x'|^4 | h_k (x', t) | & \leq   c |x'|^{-n-2} t,\\
|\frac{d}{d t}  D_{x_k}  f _k ( x', t)|,\quad |x'|^2 |\frac{d}{d t}  g _k ( x', t)|,\quad  |x'|^4 |\frac{d}{d t} h_k ( x', t)|  &  \leq c |x'|^{-n-2}.
\end{split}
\end{align}



 
 

Applying    \eqref{260318-3} and \eqref{260404-6} to  \eqref{260318-1}, if $ M'(t) < -e_0$, then  for $ 1 < t < 2$,  we obtain 
 \begin{align}\label{0915-1-1}
 \begin{split}
  \frac{d}{dt} D^2_{x_n} w_n (x', 0, t) & \leq  c|x'|^{-n-2} \int_0^t \Big( |\phi(s)| + |\phi'(s)|  + |\phi''(s)|   \Big) ( t-s)^{\frac12} ds \\
& \quad +  c|x'|^{-n-2} \int_0^t \Big( |\phi(s)| + |\phi'(s)|     \Big) ( t-s)^{-\frac12} ds+ |x'|^{-n } M'(t)\\
  & \leq   c_2 |x'|^{-n-2}   - c e_0 |x'|^{-n}\\
  & < 0
 \end{split}
 \end{align}
provided $ \frac{c_2}{c e_0}   < |x'|^2 $.
This implies  that $ D^2_{x_n}  w_n (x', 0, t) $ is strictly decreasing for $ t \in (1,2)$. This completes the proof of (ii) of 
Theorem \ref{theorem1221-2}.
\end{proof}

\end{document}
